\section{What the gain algebra establishes}
\label{sec:audited-algebra}

\paragraph{Pair scaling and its causal limit.}
For fixed content vectors, the rotary-position logit~\citep{su2024roformer}
can be written as a finite trigonometric sum
\begin{equation}
 \ell(\Delta)=\frac{1}{\sqrt{d_k}}
 \sum_j \operatorname{Re}\bigl[c_j e^{i\theta_j\Delta}\bigr].
\end{equation}
Applying a shared multiplier $m_j$ to the query and key coordinates of pair $j$
multiplies that term by $m_j^2$. Its sign and phase still depend on $c_j$.
Increasing a high-frequency coefficient therefore does not by itself establish
more local attention, higher image-band power, or sharper decoded detail.
The three existing Krea development probes have inconsistent response signs:
one prompt has no positive decoded-width diagonal among six tested bands.
Those probes jointly perturb both spatial axes on NTK2, so they also cannot
identify a single-axis response under released SEGA.

\paragraph{Standardized profile shape.}
Let $\ell$ be a nonconstant vector of sampled log powers and let
$z(\ell)=(\ell-\operatorname{mean}\ell)/\operatorname{std}(\ell)$,
using one fixed standard-deviation convention. For $\beta>0$,
\begin{equation}
 z(a\mathbf{1}+\beta\ell)=z(\ell).
\end{equation}
This elementary identity applies to an ideal standardized-profile allocation.
It does not imply invariance of SEGA's radial spread or its complete gain.
The paper-inspired reference implementation adds a standard-deviation epsilon,
and the released implementation floors energy and standard deviation. Their
behavior is therefore invariant only under the corresponding inactive-floor
or negligible-stabilizer conditions, up to floating-point error. The statement
is about a formed log-power profile; raising a two-dimensional spectrum to a
power need not commute with averaging it into that profile.

\paragraph{Allocation and scalar magnitude are coupled.}
Before clipping, write a centered SEGA multiplier as
$m_j=m_{\mathrm{ref}}(1+a_j)$ with $\operatorname{mean}a_j=0$. Then
\begin{equation}
 \operatorname{mean}(m_j^2)
 =m_{\mathrm{ref}}^2\bigl[1+\operatorname{mean}(a_j^2)\bigr].
\end{equation}
Thus centering query/key multipliers does not preserve their mean squared
magnitude, which is the mean coefficient gain. Clipping changes these moments
again. The new RMS controls isolate this effect per spatial axis. They do not
match the actual attention-logit distribution, its variance, or its entropy.

\paragraph{Forward-noise expectation and the gate.}
For the linear data--noise interpolation used to motivate rectified
flow~\citep{liu2023flow}, assume independent data $x_0$ and unit white noise
$\varepsilon$, with $x_\sigma=(1-\sigma)x_0+\sigma\varepsilon$.
With an orthonormal transform, channel summation over $C$ channels, and
mean-per-band density on non-DC bands,
\begin{equation}
 \mathbb{E}_{\varepsilon} P_b(x_\sigma)
 =(1-\sigma)^2P_b(x_0)+\sigma^2C.
\end{equation}
The cross term vanishes in expectation. This is not an exact decomposition of
one noisy sample or one generated ODE trajectory. The controller's statistic
$\rho_b=\max(P_b(x_\sigma)/(\sigma^2C)-1,0)$ and gate
$w_b=\rho_b/(1+\rho_b)$ are design choices inspired by this identity. They do
not estimate the uncertainty of the delayed clean prediction. Likewise, zero
correction at an anchored target spectrum is an algebraic property of the
rule, not a stability theorem for the full sampling process.
